\SourceLeaf{361}

\Needspace{13\baselineskip}
\begin{center}
[IV.]
\end{center}

\begin{center}
\bfseries STEREOGRAFISCHE PROJEKTION DER KUGELOBERFLÄCHE.
\end{center}

Konjugierte komplexe Zahlen sollen hier stets durch denselben Buchstaben in Klein- und Großschreibung bezeichnet werden.

Der Radius der Kugel wird $=1$ gesetzt.

$n$: lineares Vergrößerungsverhältnis.

1. Für den durch die komplexe Zahl $a$ bezeichneten Punkt gilt
\[
n = 1+aA.
\]

2. Der gegenüberliegende Punkt wird durch die komplexe Zahl $-\dfrac{1}{A}$ bezeichnet. Für diesen gilt also
\[
n = 1+\frac{1}{aA} = \frac{1+aA}{aA}.
\]

3. Der bikomplexen Zahl $a$ entspricht die trikomplexe Zahl
\[
\frac{2a+(1-aA)\eta}{1+aA}.
\]

4. Allgemeine Transformationsformel:
\[
t' = \frac{at-b}{ct+d},
\]

wobei sich $t$, $t'$ auf ein und denselben Punkt in der ersten beziehungsweise zweiten Darstellungsform beziehen.

Hier gilt
\[
n' = \text{Mod.}\ \frac{ad+bc}{(ct+d)^2}(1+tT).
\]

\Needspace{11\baselineskip}Der kleinste Wert wird erreicht für
\[
t = \frac{C}{D}
\]

\noindent und beträgt
\[
\text{Mod.}\ \frac{ad+bc}{cC+dD}\cdot\frac{D}{d} = \text{Mod.}\ \frac{ad+bc}{cC+dD}.
\]
